\section{引言}

本次作业要求数出照片里有多少粒米，并解决连体米粒计数不准的问题。

恰好我大一大二打过两次电赛，负责视觉循迹，目标识别。于是我先利用已有的知识进行了初步尝试。

我先按最基本的思路做了一遍。先转灰度，再用 Otsu 法\cite{otsu1979}自动选一个阈值把米粒和背景分开，最后统计连通域的个数。连通域指二值图中连在一起的一片前景像素，正常情况下一粒米就是一个连通域。米粒互不接触时，这一流程没有问题。一旦两粒米挨在一起，它们在二值图中连成一块，只被计为 1 粒。米粒越密，漏计越多，见图 \ref{fig:problem}。
实测粘连率 80\% 时平均少数 38.9 粒，这条误差曲线见图 \ref{fig:curve}，也正是本次作业的难点所在。

\begin{cfigure}
  \centering
  \includegraphics[width=\textwidth]{problem_zoom.png}
  \caption{粘连导致漏计的局部放大。右图中橙色的一整块是一个连通域，其中实际有 2 粒米。}
  \label{fig:problem}
\end{cfigure}

于是我转向网上学习新方法，在查阅文献时注意到，这个问题在已发表的工作中也常被回避。
Kurade 等人\cite{kurade2023}的米质检测系统采用“腐蚀取标记结合注水分割”，
注水分割即 watershed 算法，常译作“分水岭算法”。
该文承认这一算法分不开互相接触的米粒，因而把这些米粒从数据集中剔除。
本文试着解决这一问题。

本文的工作可以概括为以下四点。
\begin{enumerate}
  \item 由程序从图像自身量出单粒米的尺寸与形状，后续判据均以此为单位，像素常数是不固定的；
  \item 不预先指定颜色通道与阈值方式，而是把若干组合逐一试过，按一个打分函数自动选取；
  \item 以“又大又凹”的形状为条件，判断连通域是否为粘连块，再用注水分割切开；
  \item 切分后按面积与形状复核，并剔除硬币、边框等非米目标。
\end{enumerate}

在三组性质不同的数据上，上述方法的平均绝对误差为 4.84、0.25 与 1.90 粒，
优于本文实现的五种经典对比方法；粘连率 80\% 时误差约为直接统计连通域的四十四分之一。
在此之外，本文还用 SAM 3 的输出作伪标签蒸馏了一个 27 万参数的小网络，
详见第 \ref{sec:student} 节，其总体误差进一步降到 0.68 粒。
代码、数据来源与复现方式见附录 \ref{app:repro}。

\section{相关工作}

分开粘连目标大致有两种思路。

一种是从轮廓入手。两粒米的接触处在轮廓上是一个凹点，找到凹点再连线切开。
Tan 等人\cite{tan2019}用这种方法处理粘连稻谷，
Avzalov 等人\cite{avzalov2025}进一步用椭圆拟合滤除假凹点。本文将后者实现为对比方法 B5。

另一种是从区域入手，先用注水分割\cite{vincent1991}切开，再修正结果。
Gamarra 等人\cite{gamarra2019}先以大小和凸性判断哪些区域需要切分，切开后再合并过碎的部分。
本文方法属于后一类，结构上与之相近；我实现该方法以后查文献，才发现有前辈用过类似的方法。

国内在谷物考种一类的任务上做过不少同类工作。凌云等人\cite{ling2005}较早把注水分割用在谷物籽粒上，
杨蜀秦等人\cite{yang2009}处理的正是连接在一起的米粒。
荀一等人\cite{xun2010}在连续腐蚀与膨胀的过程中找出两粒之间的公共区域，取其端点连线切开。
柳冠伊等人\cite{liu2014}把距离变换图当作地形图，沿山谷线切分。
他们还指出，这样一来就不必先由人工设定阈值去判断是否粘连。
这一点与本文用自标定量出判据的想法相近。
李文勇等人\cite{li2015}则先用形状因子判断一个区域是不是粘连，再定位分割点。
判断的那一步与本文第 \ref{sec:touch} 节做的是同一件事，只是所用的量不同。

\section{方法一：几何方法}

方法一沿用常规的思路，先二值化，再找连通域，最后用注水分割把粘连的米粒切开。
不同之处在于，判断时用到的尺寸和门限都不预先设定，而是从每张图自身量出来。

整个流程分五步，见图 \ref{fig:flow}。第一步选出米粒所在的前景，第二步量出一粒米有多大，
第三步找出哪些连通域是几粒米粘在一起，第四步把这些粘连块切开，
第五步复核切分结果，并剔除硬币、木框这类异物。各步在一张合成图上的中间结果见图 \ref{fig:pipeline}。
下面先交代几个后文要用的概念，再按这五步依次说明。

\begin{cfigure}
  \centering
  % !TEX root = 米粒粘连导致漏数的问题与解决.tex
% 方法流程图。左列为整图处理，右列为对每个连通域的判断与计数。
% 节点只设最小宽度、不设 text width，避免中文在定宽盒子里被两端对齐拉开字距。
\begin{tikzpicture}[
  font=\small,
  node distance=5mm and 7mm,
  box/.style={draw, rounded corners=2pt, align=center, fill=blue!6,
              minimum height=9mm, minimum width=44mm, inner sep=3pt},
  dec/.style={draw, diamond, aspect=2.8, align=center, fill=orange!10,
              inner sep=1pt, minimum width=36mm},
  side/.style={draw, rounded corners=2pt, align=center, fill=gray!10,
               minimum height=8mm, minimum width=24mm, inner sep=3pt},
  io/.style={box, fill=green!10},
  arr/.style={-{Stealth[length=2mm]}, thick},
  lab/.style={font=\footnotesize, inner sep=1.5pt},
]
  % ---- 左列：整图处理 ----
  \node[io]  (in)    {输入图像};
  \node[box, below=of in]    (bin)   {多候选二值化\\通道 × 类数 × 前景亮暗};
  \node[box, below=of bin]   (score) {按 $\mathrm{score}=N\sub{单}A_0\gamma$\\选出得分最高的候选};
  \node[dec, below=of score] (chk)   {合理性检查\\通过？};
  \node[box, below=7mm of chk] (cal) {尺度自标定\\量出 $A_0$、$b_0$、$\rho_0$、$w_0$};

  \draw[arr] (in) -- (bin);
  \draw[arr] (bin) -- (score);
  \draw[arr] (score) -- (chk);
  \draw[arr] (chk) -- node[lab, right] {是} (cal);
  \draw[arr] (chk.west) -- node[lab, above] {否} ++(-6mm,0)
             |- (score.west);
  \node[lab, rotate=90, anchor=south]
       at ($($(chk.west)+(-6mm,0)$)!0.5!($(chk.west |- score.west)+(-6mm,0)$)$) {换下一候选};

  % ---- 右列：逐个连通域 ----
  \node[box, right=29.4mm of in]  (comp) {逐个取出连通域};
  \node[dec, below=of comp]     (d1)   {面积 $<0.45A_0$？};
  \node[dec, below=of d1]       (d2)   {异物、亮边\\或暗区？};
  \node[dec, below=of d2]       (d3)   {是粘连块？};
  \node[box, below=of d3]       (ws)   {h-maxima 放种子\\注水分割切开};
  \node[box, below=of ws]       (fix)  {碎块合并；仍大且凹则\\重切，不成按面积补数};
  \node[io, below=of fix]       (sum)  {累加，再从开运算前的\\阈值图找回被抹掉的米};

  \node[side, right=of d1] (s1) {碎屑，忽略};
  \node[side, right=of d2] (s2) {不计数};
  \node[side, right=of d3] (s3) {计 1 粒};

  \draw[arr] (cal.east) -- ++(8mm,0) |- (comp.west);
  \draw[arr] (comp) -- (d1);
  \draw[arr] (d1) -- node[lab, right] {否} (d2);
  \draw[arr] (d2) -- node[lab, right] {否} (d3);
  \draw[arr] (d3) -- node[lab, right] {是} (ws);
  \draw[arr] (ws) -- (fix);
  \draw[arr] (fix) -- (sum);
  \draw[arr] (d1) -- node[lab, above] {是} (s1);
  \draw[arr] (d2) -- node[lab, above] {是} (s2);
  \draw[arr] (d3) -- node[lab, above] {否} (s3);
  \draw[arr] (s3) |- (sum.east);
\end{tikzpicture}

  \caption{方法流程图。}
  \label{fig:flow}
\end{cfigure}

\begin{cfigure}
  \centering
  \includegraphics[width=\textwidth]{pipeline_synthetic.png}
  \caption{各步骤在一张合成图上的中间结果。}
  \label{fig:pipeline}
\end{cfigure}

\subsection{几个基本概念}

\textbf{距离变换 $D$} 是每个像素到最近背景像素的距离。
米粒中间的值大、边缘的值小，像一座山；两粒米贴合的地方是一道山谷。

\textbf{种子} 是告诉程序“这里有一粒米”的起点，只需落在米粒内部的一小块像素上。

\textbf{注水分割} 把距离变换取负值当作地形，每粒米陷成一个盆地，见图 \ref{fig:terrain}。
水从各个种子处涨起来，两股水相遇的地方就是切分线。
给几个种子就切成几块，所以一个粘连块被算作几粒，取决于种子怎么放。

\begin{cfigure}
  \centering
  \includegraphics[width=\textwidth]{terrain_3d.png}
  \caption{把距离变换取负值当作地形。(a) 为一个两粒粘连块，(b) 为它的三维地形，
  竖轴为距背景的像素距离，(c) 为注水后的切分结果。}
  \label{fig:terrain}
\end{cfigure}

\textbf{凸实度 $\rho$} 是区域面积与其凸包面积之比，凸包即包住该区域的最小凸多边形。
一粒米接近凸形，$\rho$ 接近 1；两粒粘在一起，接缝处凹进去，$\rho$ 会变小。

\subsection{第一步，选出米粒所在的前景}

最简单的做法是用 Otsu 法把灰度图分成亮暗两类。可是真实照片里，白米放在深色衬布上，
衬布又铺在浅色桌面上，画面有三个亮度层次。分两类只能切一刀，
无论切在哪里，总有一整片背景被划进前景，见图 \ref{fig:binfail}。
这片伪前景比所有米粒加起来还大，后面每一步都会被它带偏。

\begin{cfigure}
  \centering
  \includegraphics[width=\textwidth]{binarisation_failure.png}
  \caption{固定用灰度两类阈值的后果，蓝色为该二值图中最大的连通域。
  前景取亮或取暗都会多出一片巨大的伪前景。}
  \label{fig:binfail}
\end{cfigure}

所以本文不事先决定用哪个颜色通道、分几类、前景取亮还是取暗，
而是把这些组合都试一遍，用下式打分，选分数最高的一个。
\begin{equation}
  \mathrm{score} = N\sub{单} \cdot A_0 \cdot \gamma
  \label{eq:score}
\end{equation}
式中 $N\sub{单}$ 是大小接近一粒米的连通域个数，$A_0$ 是估计出的单粒面积，
$\gamma$ 是前景与紧邻背景的反差。前两项奖励“分出了许多大小一致的米粒”；
第三项用来排除噪声斑点，它们数量很多，却几乎没有反差。

选中的方案还要过一道合理性检查，米粒的长宽比和大小要在合理范围之内，
不通过就换下一个候选。低分辨率图里被误当成米粒的木框条纹，就是靠这道检查挡下的。

\subsection{第二步，量出一粒米有多大}

后面判断多大算大、多凹算凹，都需要一粒米作参照。
三组数据的米粒大小相差好几倍，固定的像素数难以通用，只能每张图现量。

做法是统计所有连通域的面积，取出现次数最多的那个大小作为单粒面积 $A_0$，
实现见代码 \ref{code:a0}。单独的米粒通常最多，粘连块的面积则落在 $A_0$ 的两倍、三倍附近。
面积先取对数，这些峰就变得等间距，更容易分辨，见图 \ref{fig:calib}。
起初取的是面积最小的那个峰，结果常常选中比米粒还小的碎屑。

有了 $A_0$，再从大小接近一粒米的连通域里量出短轴 $b_0$、凸实度 $\rho_0$ 和宽度 $w_0$，
作为后面各步的尺子。

\begin{lstlisting}[float=!htbp, caption={估计单粒面积 $A_0$}, label={code:a0}]
def estimate_a0(areas):
    """取对数面积分布的最高峰作为单粒米面积"""
    log_s = np.log(areas)
    kde = gaussian_kde(log_s)                 # 核密度估计，相当于平滑的直方图
    grid = np.linspace(log_s.min(), log_s.max(), 512)
    a0 = np.exp(grid[np.argmax(kde(grid))])   # 最高峰对应的面积
    # 再取 0.65~1.45 倍 A0 之间的连通域求平均，结果更稳
    band = areas[(areas >= 0.65 * a0) & (areas <= 1.45 * a0)]
    return band.mean() if len(band) > 0 else a0
\end{lstlisting}

\begin{cfigure}
  \centering
  \includegraphics[width=\textwidth]{calibration.pdf}
  \caption{面积分布的最高峰即单粒面积 $A_0$，粘连块集中在它的整数倍附近。}
  \label{fig:calib}
\end{cfigure}

\subsection{第三步，找出粘连块}
\label{sec:touch}

比一粒米大的连通域，未必是几粒粘在一起，也可能只是一粒偏大的米。
两粒米贴合处会凹进去一块，所以本文要求又大又凹才算粘连，实现见代码 \ref{code:touch}。
起初只要大或凹占一条就算，结果在真实照片上把不少单粒也切开了。

米粒只有几个像素宽时，边缘锯齿很重，凸实度量不准，这时改看厚度。
区域里能放下的最大圆的直径记为宽度 $w$。两粒米并排时面积翻倍，宽度却几乎不变，
所以面积与宽度平方之比随粒数增大，一粒米则基本不变，见图 \ref{fig:criterion}。

\begin{cfigure}
  \centering
  \includegraphics[width=\textwidth]{touching_criterion.png}
  \caption{单粒与粘连块的对照。蓝色为区域本身，浅色为凸包多出的部分，橙圈为最大内切圆。}
  \label{fig:criterion}
\end{cfigure}

\begin{lstlisting}[float=!htbp, caption={判断一个连通域是不是粘连块}, label={code:touch}]
def is_touching(c, cal):
    """c 是一个连通域，cal 是自标定得到的单粒米参数"""
    if c.area <= 1.2 * cal.a0:            # 不比一粒米大，不算粘连
        return False
    if cal.b0 < 4.5:                      # 米粒太小，凸实度量不准，改看厚度
        e = (c.area / cal.a0) / (c.width / cal.w0) ** 2
        return e > 1.5
    return c.solidity < cal.rho0 - 0.06   # 又大又凹才算粘连
\end{lstlisting}

\subsection{第四步，切开粘连块}

判出粘连块之后，要决定放几个种子，也就是认定块里有几粒米。
直接限制种子之间的距离行不通。米粒细长，一粒米内部的山脊比两粒并排时的间距还长，
距离定多少都顾此失彼。

关键在于看深度而不是距离。两粒米之间是一道深谷，一粒米内部只有浅浅的起伏。
于是对距离变换做 h-maxima，只保留比周围高出 $h$ 以上的峰作种子，
见图 \ref{fig:seeddepth} 与代码 \ref{code:split}。
$h$ 随短轴 $b_0$ 缩放，因为两粒米能靠多近取决于它们的宽度；
但不低于一个像素，因为距离变换本身就以像素计量，更小的起伏只是边缘的锯齿。

\begin{cfigure}
  \centering
  \includegraphics[width=\textwidth]{seed_depth.pdf}
  \caption{为什么看深度而不看距离。左为粘连块的距离变换，右为沿山脊的剖面。
  谷底低于橙色虚线，即比较低的峰矮出 $h$ 以上，两个峰各留一个种子。}
  \label{fig:seeddepth}
\end{cfigure}

这一步有个问题，是在检查效果图时才发现的。h-maxima 作用在浮点的距离图上，
会把同一个平顶的峰碎成几小块，每块都成了种子，一粒米就被切成两半。
同一粒米上的碎块彼此很近，不同米粒的种子则隔得较远，
所以把相距不到约半粒米宽的种子并成一颗。

这个问题还连着另一个。碎块没合并时，$h$ 越小碎块越多，调参只好把 $h$ 调大；
$h$ 太大，又把饱满米粒之间那道浅颈抹平了，许多粘连块切不开。
两种错误在数量上碰巧抵消，总数看着是对的，效果图却明显不对。
合并碎块之后再把 $h$ 调小，两种错误都基本消失，量化结果见第 \ref{sec:ablation} 节。

\begin{lstlisting}[float=!htbp, caption={放种子并做注水分割}, label={code:split}]
def split_cluster(mask, cal, beta=0.1):
    dist = cv2.distanceTransform(mask, cv2.DIST_L2, 5)
    h = max(beta * cal.b0 / 2, 1.0)        # 深度门限跟着米粒宽度变，不低于 1 像素
    peaks = h_maxima(dist, h)              # 只留下足够深的山峰
    r = round(0.6 * cal.b0 / 2)            # 相距不到 0.6 倍短轴的峰算同一粒米
    seeds = label(cv2.dilate(peaks, disk(r))) * peaks
    if seeds.max() <= 1:                   # 只有一个种子，不用切
        return (mask > 0).astype(int)
    return watershed(-dist, seeds, mask=mask > 0)
\end{lstlisting}

\subsection{第五步，复核并剔除异物}
\label{sec:correct}

注水分割只给出了可能的切法，切得对不对还要复核。
一个粘连块从切开到定出粒数的全过程见图 \ref{fig:cluster}，计数过程见代码 \ref{code:count}。
复核做三件事。

一是并回碎块。比半粒米还小的碎块多半是切多了，把它并回与它接壤最多的那一块。

二是补切。切完后仍然又大又凹的块，说明没有切开。先按面积算出它该有几粒，
再取距离图上最高的那几个峰作种子重切一次。切出的每一块大小都像一粒米才采用，
否则只按面积计数。起初这里只按面积计数，粒数大多对，但效果图上几粒米画成一整块。

三是剔除异物。硬币、木框这类东西面积很大，但与米堆有两点不同。
米粒只是挨着，不会融成一体，所以米堆最厚处也只有一粒米那么厚，硬币和木框则厚得多。
几粒米粘连通常会留下凹口，面积数倍于单粒却和单粒一样凸的区域，多半不是米堆。
这条判据起初还要求异物比米粒更圆，细长的木框条因此漏了过去，去掉这一条后才被挡住。

\begin{lstlisting}[float=!htbp, caption={计数主循环}, label={code:count}]
total = 0
for c in components:
    if c.area < 0.3 * a0:                  # 碎屑，忽略
        continue
    if is_foreign(c, cal):                 # 硬币、木框等，不计数
        continue
    if not is_touching(c, cal):            # 单独的一粒
        total += 1
        continue
    parts = merge_small(split_cluster(c.mask, cal), a0)  # 小于 0.4*A0 的碎块并入邻居
    for p in regionprops(parts):
        k = p.area / a0
        if k < 0.4:                        # 没有邻居可并的碎块
            continue
        if k > 1.5 and p.solidity < cal.rho0 - 0.03:
            n = round(k)                   # 没切开，按面积它该有 n 粒
            pieces = split_by_count(p.mask, n, cal)   # 取最高的 n 个峰重切一次
            if pieces_look_like_grains(pieces, a0):
                parts_drawn.extend(pieces)            # 切对了就按切开的画
            total += n                     # 粒数都按 n 计
        else:
            total += 1
\end{lstlisting}

\begin{cfigure}
  \centering
  \includegraphics[width=\textwidth]{cluster_detail.png}
  \caption{粘连块的处理过程，自左至右为粘连块、距离变换、种子与最终结果。}
  \label{fig:cluster}
\end{cfigure}

\section{方法一的实验与分析}

这一节用实验回答五个问题。方法一比常见做法准多少，粘连加重时还准不准，
图像变差时从哪里开始失效，每一步各起多大作用，以及和视觉大模型 SAM 3 相比差在哪里、省在哪里。

\subsection{数据集}

三组数据中，D1、D3 与 D2 所用的米粒素材取自 Roboflow Universe 上的公开数据集，许可为 CC BY 4.0。
第 \ref{sec:student} 节的方法二沿用这三组数据。

\textbf{D1，真实照片，51 张。} 取自 Roboflow 上的 rice 数据集\cite{ds_d1}，
手机拍摄，深色布上摊开的米粒，每张 80--116 粒，带一枚参照硬币，每粒米都有人工标注框。
米粒基本不粘连。

\textbf{D2，自制合成图，40 张。} 从 RICE 数据集\cite{ds_grains}中 Karacadag 品种的单粒米照片里
抠出米粒，随机旋转、缩放后贴到深色画布上，
粘连率分 0\%、20\%、40\%、60\%、80\% 五档，每档 8 张，每张 44--98 粒。
它的真值由构造方式精确给定，粘连程度也可控，专门用来考粘连。

\textbf{D3，低分辨率，719 张。} 取自 Roboflow 上的 RICE 数据集\cite{ds_d3}，
$224\times224$ 的小图，每张约 14 粒。该数据集导出时做过翻转增强，同一源图存在多个副本，
按源图名去重后保留 719 张。

\subsection{评价指标}

记第 $k$ 张图的真值为 $g_k$、预测为 $p_k$。
主要指标为平均绝对误差 $\mathrm{MAE}=\frac1n\sum_k|p_k-g_k|$，即平均每张数错几粒。
另外给出准确率 $1-\mathrm{MAPE}$ 与最差单张误差，
其中 MAPE 为平均相对误差 $\frac1n\sum_k|p_k-g_k|/g_k$。
平均值小并不排除个别图像错得离谱，所以两者一并看。

\subsection{与其他方法比较}

先看方法一比常见做法准多少。
为使比较只落在“如何分开粘连”这一步上，所有方法都跑在同一份二值图上。
B1 直接统计连通域个数，不对粘连作任何处理。
B2 用前景总面积除以单粒面积。
B3 是距离变换注水分割，种子取 $D > 0.5\max D$，这是 OpenCV 教程的做法。
B4 腐蚀后取连通域作种子，即文献\cite{kurade2023}的做法。
B5 为轮廓凹点 + 椭圆拟合\cite{avzalov2025}。

结果见表 \ref{tab:comparison}。三组数据上，方法一的平均误差都是六种几何方法中最低的。
三种对比方法的失效值得单独说一下。

\begin{ctable}
  \centering
  \caption{三组数据上的总体结果，误差以粒计，“最差”一列为最差单张误差。
  加粗为六种几何方法中每列最好的值，SAM 3 另列在最后作参照。
  B5 和 SAM 3 是在每个数据集上分别调到最好后的结果。}
  \label{tab:comparison}
  \tabfit{\setlength{\tabcolsep}{4pt}
\begin{tabular}{lrrrrrrrrr}
\toprule
& \multicolumn{3}{c}{D1 真实稀疏} & \multicolumn{3}{c}{D2 合成粘连} & \multicolumn{3}{c}{D3 低分辨率} \\
\cmidrule(lr){2-4}\cmidrule(lr){5-7}\cmidrule(lr){8-10}
方法 & MAE & 准确率\% & 最差 & MAE & 准确率\% & 最差 & MAE & 准确率\% & 最差 \\
\midrule
B1 连通域计数 & 6.31 & 93.25 & 19 & 22.95 & 67.94 & 51 & 3.31 & 76.17 & \textbf{16} \\
B2 面积估算 & 22.76 & 74.72 & 87 & 2.30 & 96.69 & 18 & 18.39 & -30.64 & 76 \\
B3 距离变换注水分割 & 74.90 & 17.77 & 123 & 8.20 & 88.63 & 24 & 8.27 & 42.82 & 18 \\
B4 腐蚀标记注水分割 \cite{kurade2023} & 5.45 & 94.21 & 26 & 22.93 & 67.97 & 51 & 4.29 & 69.37 & \textbf{16} \\
B5 凹点+椭圆拟合 \cite{avzalov2025} & 5.76 & -- & -- & 2.08 & -- & -- & 3.32 & -- & -- \\
\textbf{本文方法} & \textbf{2.75} & \textbf{96.97} & \textbf{18} & \textbf{0.25} & \textbf{99.61} & \textbf{3} & \textbf{0.50} & \textbf{96.11} & \textbf{16} \\
\midrule
SAM 3 零样本 & 3.61 & 96.11 & 24 & 72.15 & 0.00 & 98 & 0.18 & 98.39 & 19 \\
\bottomrule
\end{tabular}}
\end{ctable}

B4 在合成数据上的误差与 B1 逐档相同。一次 $3\times3$ 腐蚀不足以分开粘连块，
每块仍只产生一个种子，注水分割什么也没切开，等于退化成连通域计数。

B3 在真实照片上误差达 74.90 粒。它的阈值取自整张图距离变换的最大值，
而画面里的硬币远大于米粒，把阈值整体抬高，大多数米粒的种子被抹掉了。
图 \ref{fig:b3seeds} 是同一张照片上两种取种子方式的对照，
B3 全图只留下 19 个种子，方法一按连通域分别标定，得到 150 个。

\begin{cfigure}
  \centering
  \includegraphics[width=\textwidth]{b3_seeds.png}
  \caption{B3 的全局阈值与本文自适应种子的对照，红色为种子，灰色为二值图。
  本文对每个连通域分别放种子，粘连块会得到不止一个，因此总数多于米粒数。}
  \label{fig:b3seeds}
\end{cfigure}

B5 难以同时兼顾两种尺度。它在轮廓上作一个小圆找凹点，米粒只有 4 像素宽时，
圆几乎盖住整粒米，孤立的单粒也会被判出凹点。把圆调大，它在 D2 上又几乎不再切分。
需要说明，原文的椭圆拟合为穷举搜索，本文实现的是贪心近似，对它未必公平。

这三种方法失效的原因相同，都用了预先设定的固定数值。
换一种米粒尺度，这些数字就不再合适，方法一的判据从图像自身量出，所以在三组数据上表现都比较稳。

\subsection{误差随粘连程度的变化}

再看粘连加重时方法一还准不准。结果见表 \ref{tab:touching} 与图 \ref{fig:curve}。
粘连率自 0\% 升至 80\%，连通域计数的误差由 0.25 粒一路升到 38.88 粒，
方法一在五档上均未超过 0.9 粒。
图 \ref{fig:scatter} 是逐张的预测值与真值对照，格子集中在对角线附近，没有整体偏多或偏少。
可见粘连加重时，方法一的误差只是缓慢上升，没有随粘连一起涨上去。

\begin{ctable}
  \centering
  \caption{合成数据上按粘连率分层的平均绝对误差，单位为粒，加粗为每一档误差最小的方法。}
  \label{tab:touching}
  \tabfit{\begin{tabular}{lrrrrrr}
\toprule
粘连率 & B1 & B2 & B3 & B4 & B5 & \textbf{本文} \\
\midrule
0\% & 0.25 & 1.12 & 0.00 & 0.12 & 0.00 & 0.25 \\
20\% & 12.50 & 1.12 & 4.50 & 12.50 & 7.62 & 0.75 \\
40\% & 27.62 & 2.25 & 10.62 & 27.62 & 15.88 & 2.38 \\
60\% & 35.50 & 2.50 & 11.12 & 35.50 & 16.38 & 2.00 \\
80\% & 38.88 & 4.50 & 14.75 & 38.88 & 18.25 & 3.50 \\
\bottomrule
\end{tabular}}
\end{ctable}

\begin{cfigure}
  \centering
  \includegraphics[width=\textwidth]{error_vs_touching.pdf}
  \caption{计数误差随粘连率的变化。}
  \label{fig:curve}
\end{cfigure}

\begin{cfigure}
  \centering
  \includegraphics[width=\textwidth]{pred_vs_true.pdf}
  \caption{预测值与真值对照的热力图。格子颜色表示落在该格的图片张数，虚线表示两者相等。}
  \label{fig:scatter}
\end{cfigure}

\subsection{图像质量变差时的表现}
\label{sec:robust}

这一小节找方法一的失效边界。本文在 D2 上人为施加模糊、噪声、对比度压缩与降采样四类退化，
结果见表 \ref{tab:robust}，其中模糊强度以单粒宽度 $b_0$ 为单位。

噪声加到 45 个灰度级、对比度压到 8\%，误差都没有上升。
降采样到 1/6 时米粒只剩约 2 像素宽，中位误差为 4 粒，误差逐步上升而不是突然失效。
模糊则不同。中位误差变化不大，平均误差却突然暴涨，
原因是个别图像模糊后二值化整个选错，一张图数出上千粒。
所以方法一怕的是模糊图像。

\begin{ctable}
  \centering
  \caption{合成粘连图 D2 在不同退化下的计数误差，共 40 张。“崩溃”指单张误差超过 100 粒，
  出现崩溃的格子加粗。}
  \label{tab:robust}
  \tabfit{\begin{tabular}{llrrrr}
\toprule
退化类型 & 强度 & 平均误差 & 中位误差 & 崩溃场景 & B1 平均误差 \\
\midrule
高斯模糊 & 0.00$\,b_0$ & 1.77 & 1.00 & 0/40 & 22.95 \\
 & 0.20$\,b_0$ & 251.88 & 1.00 & 2/40 & 156.95 \\
 & 0.35$\,b_0$ & 79.28 & 2.00 & 2/40 & 65.78 \\
 & 0.70$\,b_0$ & 77.35 & 13.50 & 13/40 & 59.95 \\
\midrule
加性噪声 & 0 & 1.77 & 1.00 & 0/40 & 22.95 \\
 & 20 & 1.60 & 1.00 & 0/40 & 23.00 \\
 & 45 & 1.57 & 1.00 & 0/40 & 22.95 \\
\midrule
对比度保留 & 8\% & 1.62 & 1.00 & 0/40 & 22.98 \\
 & 30\% & 1.85 & 1.00 & 0/40 & 22.95 \\
 & 100\% & 1.77 & 1.00 & 0/40 & 22.95 \\
\midrule
降采样倍率 & 1.0$\times$ & 1.77 & 1.00 & 0/40 & 22.95 \\
 & 3.0$\times$ & 2.30 & 1.00 & 0/40 & 23.88 \\
 & 6.0$\times$ & 5.35 & 4.00 & 0/40 & 26.00 \\
\bottomrule
\end{tabular}}
\end{ctable}

\subsection{消融实验}
\label{sec:ablation}

这一小节逐项关掉或改动方法一的某一步，看每一步各起多大作用，结果见表 \ref{tab:ablation}。

\begin{ctable}
  \centering
  \caption{去掉或改动某一步后的平均绝对误差，单位为粒，第一行为完整方法。
  每列加粗的是误差最大的一项，即去掉后损失最大的那一步。}
  \label{tab:ablation}
  \tabfit{\begin{tabular}{lrrr}
\toprule
消融项 & D1 & D2 & D3 \\
\midrule
完整方法 & 2.75 & 0.25 & 0.50 \\
去掉尺度自标定，$A_0$ 取数据集均值 & \textbf{6.31} & 0.17 & 0.47 \\
去掉形状先验校正 & 4.45 & 1.12 & 0.42 \\
去掉异物剔除 & 4.25 & 0.25 & \textbf{2.00} \\
异物剔除去掉宽度判据 & 3.90 & 0.25 & 0.66 \\
异物判据附加“更圆”条件，即最初写法 & 3.90 & 0.25 & 0.98 \\
米粒短轴不足 4.5 像素时改用厚度判据判粘连，即修正前的写法 & 4.51 & 0.25 & 0.38 \\
去掉亮背景判据 & 3.25 & 0.25 & 0.45 \\
碎屑门限取 0.3，即修正前的写法 & 2.78 & 0.25 & 0.66 \\
去掉暗区判据 & 3.53 & 0.25 & 0.90 \\
不找回被开运算抹掉的米粒 & 4.22 & 0.25 & 0.74 \\
粘连判据改为“或” & 2.90 & 0.07 & 0.74 \\
种子碎块不合并 & 2.76 & \textbf{3.17} & 0.66 \\
种子碎块不合并，$\beta$ 取 0.45，即修正前的写法 & 2.65 & 1.77 & 0.66 \\
仅灰度通道 & 2.33 & 0.25 & 1.00 \\
加 $3\times3$ 中值滤波 & 5.47 & 0.23 & -- \\
\bottomrule
\end{tabular}}
\end{ctable}

合并种子碎块管的是切分本身。不合并时，D2 的误差从 0.25 粒涨到 3.17 粒，
一粒米常被切成两半。修正前的写法不合并碎块，只好把 $\beta$ 抬到 0.45 来压住碎块，
代价是许多粘连块切不开，D2 的误差为 1.77 粒。

形状先验校正，也就是按面积补数，兜底的是注水分割仍没切开的那些块。
去掉它，D2 的误差从 0.25 粒涨到 1.12 粒，D1 从 4.84 粒涨到 5.31 粒。

异物剔除管的是 D3 的木框。去掉它，D3 的误差从 1.90 粒涨到 3.74 粒。

尺度自标定管的是真实照片里参差的米粒大小。改成整个数据集共用一个 $A_0$ 后，
D1 的误差从 4.84 粒涨到 8.20 粒。D2 反而从 0.25 粒降到 0.17 粒，
因为合成图的米粒本来就大小整齐，一个固定值比每张图现量更稳。

只用灰度通道时 D1 稍好，为 4.63 粒，但 D3 离不开饱和度通道，所以候选通道仍都保留。

表中也有两项在个别数据上反而更好。粘连判据改成“或”后 D2、D3 的误差略降，
去掉形状先验校正后 D3 也略降，但它们都让真实照片 D1 变差，所以没有采用。

这几步各管一类问题。去掉其中哪一步，都有一组数据明显变差。

\subsection{与 SAM 3 对比}
\label{sec:sam3}

最后把方法一和视觉大模型 SAM 3\cite{sam3} 放在一起比。
SAM 3 是去年冬天由 Meta 推出的视觉分割模型，由文本提示驱动，
给出“rice grain”之类的提示，就能返回每粒米的掩膜，无需训练。
本文先在 D1 上比较了不同提示词与阈值，结果见表 \ref{tab:sam3}，
再把各提示词在三组数据上分别调到最好，结果见表 \ref{tab:sam3matrix}。

\begin{ctable}
  \centering
  \caption{SAM 3 在真实照片 D1 上使用不同提示词和阈值的平均绝对误差，单位为粒，加粗为误差最小的一格。}
  \label{tab:sam3}
  \tabfit{\begin{tabular}{lrrrrrr}
\toprule
文本提示 & $\tau=0.15$ & $\tau=0.20$ & $\tau=0.25$ & $\tau=0.30$ & $\tau=0.35$ & $\tau=0.40$ \\
\midrule
\texttt{grain} & 18.86 & 9.04 & 4.69 & 3.63 & 3.61 & 4.53 \\
\texttt{rice grain} & 8.71 & 6.41 & 8.12 & 13.45 & 28.41 & 51.78 \\
\texttt{white rice grain} & 18.10 & 10.00 & 7.16 & 5.92 & 5.88 & 6.71 \\
\bottomrule
\end{tabular}}
\end{ctable}

\begin{ctable}
  \centering
  \caption{各提示词在各数据集上分别调到最好时的误差，$\tau^\ast$ 为对应阈值，
  加粗为每组数据上误差最小的提示词。}
  \label{tab:sam3matrix}
  \tabfit{\begin{tabular}{lrrrrrr}
\toprule
& \multicolumn{2}{c}{D1} & \multicolumn{2}{c}{D2} & \multicolumn{2}{c}{D3} \\
\cmidrule(lr){2-3}\cmidrule(lr){4-5}\cmidrule(lr){6-7}
文本提示 & MAE & $\tau^\ast$ & MAE & $\tau^\ast$ & MAE & $\tau^\ast$ \\
\midrule
\texttt{rice grain} & 6.41 & 0.20 & 72.15 & 0.05 & 0.62 & 0.20 \\
\texttt{grain} & 3.61 & 0.35 & 72.15 & 0.05 & \textbf{0.07} & 0.35 \\
\texttt{white rice grain} & 5.88 & 0.35 & 72.12 & 0.05 & 0.08 & 0.20 \\
\texttt{white seed} & \textbf{3.29} & 0.40 & \textbf{0.65} & 0.40 & \textbf{0.07} & 0.35 \\
\texttt{seed} & 24.45 & 0.25 & 27.45 & 0.20 & 0.08 & 0.40 \\
\bottomrule
\end{tabular}}
\end{ctable}

当提示词选得恰当时，SAM 3 在 D1、D3 上比方法一更准，在重度粘连的 D2 上则是方法一更准。
\texttt{white seed} 在三组数据上分别为 3.29、0.65、0.07 粒。
但它对提示词相当敏感，五个提示词中有三个在合成图上一粒也没检出，
其中包括看起来最自然的 \texttt{rice grain}，而模型不会提示自己失败。

图 \ref{fig:render_d1} 至图 \ref{fig:render_d3} 是每组数据一个场景的逐粒结果，便于核对。
每粒米涂一种颜色，相邻两粒颜色不同。
虚线框与 $\times N$ 标出的区域并未被真正切开，其粒数由面积推断而来。
灰色并标 $\times$ 者为判定的异物，不计入总数。

\begin{cfigure}
  \includegraphics[width=\textwidth]{render_d1.png}
  \caption{真实照片 D1 上的结果，左为本文方法、右为 SAM 3。}
  \label{fig:render_d1}
\end{cfigure}

\begin{cfigure}
  \includegraphics[width=\textwidth]{render_d2.png}
  \caption{合成粘连图 D2 中重度粘连场景的结果，左为本文方法、右为 SAM 3。}
  \label{fig:render_d2}
\end{cfigure}

\begin{cfigure}
  \includegraphics[width=\textwidth]{render_d3.png}
  \caption{低分辨率图 D3 上的结果，左为本文方法、右为 SAM 3。}
  \label{fig:render_d3}
\end{cfigure}

D1 中 SAM 3 没有把硬币计入，方法一却把硬币上的浮雕碎片数成了 8 粒，详见附录 \ref{app:renders}。
D2 中方法一数出 63 粒，真值 64 粒，每个粘连块都被真正切开，没有靠面积折算的地方。
D3 中木框亮边被正确剔除，但漏掉了一粒颜色偏暗的米。

\subsection{算力成本对比}
\label{sec:cost}

准确度之外，两种方法的算力开销差别很大。
测量在同一台笔记本上进行，配置为 Intel i9-13900H 处理器与 12 GB 显存的 RTX 4080 Laptop GPU。
每组数据取 10 张图，测每张的处理时间，不含读图，并记录峰值内存和显存，结果见表 \ref{tab:cost}。
SAM 3 在 CPU 上太慢，每组只测了 1 张。CPU 上用 32 位浮点，GPU 上用 16 位浮点 bfloat16。


\begin{ctable}
  \centering
  \caption{单张图的处理时间和资源占用，显存为分配峰值，SAM 3 在 CPU 上每组只测 1 张。
  每行加粗的是最省的一项。}
  \label{tab:cost}
  \tabfit{\begin{tabular}{lrrrrr}
\toprule
项目 & 本文 CPU & 本文 GPU & 学生模型 GPU & SAM 3 GPU & SAM 3 CPU \\
\midrule
D1 每张耗时／s & 0.37 & 0.15 & 0.014 & 0.21 & 17.7 \\
D2 每张耗时／s & 0.63 & 0.40 & 0.025 & 0.21 & 17.3 \\
D3 每张耗时／s & 0.033 & 0.071 & 0.002 & 0.21 & 18.3 \\
\midrule
模型参数量 & 0 & 0 & 0.27 M & 840 M & 840 M \\
权重文件／MB & 0 & 0 & 1 & 3281 & 3281 \\
模型加载／s & -- & -- & 0.1 & 3.1 & 2.5 \\
峰值内存／GB & 0.30 & 0.70 & 1.17 & 5.61 & 5.07 \\
峰值显存／GB & -- & 0.10 & 0.23 & 2.50 & -- \\
\bottomrule
\end{tabular}}
\end{ctable}

有显卡时 SAM 3 每张约 0.21 s，方法一的 GPU 版在 D1 上更快，为 0.15 s，在 D2 上略慢。
但 SAM 3 需先加载 8.4 亿参数、3.2 GB 权重，运行时占用 5.6 GB 内存与 2.5 GB 显存，
方法一没有需要学习的参数，GPU 版只占 0.09 GB 显存。
没有显卡时差距更大。SAM 3 每张耗时 17 s 以上，方法一在 CPU 上为 0.03 至 0.63 s。
也就是说，有显卡时两者速度相近，没有显卡时方法一快得多，而且几乎不占内存。

\subsection{本节小结}

方法一在三组数据上都比五种常见做法准，粘连加重时误差也几乎不上升。
它的弱点有两处，一是怕模糊，二是在低分辨率的 D3 上与 SAM 3 差距明显。
SAM 3 在真实照片和低分辨率图上更准，却要挑提示词、要显卡。
有没有方法能能兼得两者的长处，于是我想到了第二种方法。

\section{方法二：SAM3作为教师模型训练专用小模型}
\label{sec:student}

\subsection{方法二思路的由来}

第二种方法是做完上一节的实验之后才想到的。那里留下一个通过几何方法难以解决的难题。
本文方法无论怎么调都与 SAM 3 差距明显。失败的样子还大多属于同一类，
木框亮边被当成米，颜色偏暗的米被漏掉。
追究下去，几何判据只知道米长什么样，并不知道什么是米，
既然 SAM 3 知道，能否把它知道的东西教给一个小模型？

\subsection{具体实现方式}

\textbf{制作图片标注：} 用 SAM 3 处理每张照片，取其每个掩膜的质心，得到一串点坐标作为伪标签。
合成图 D2 则直接采用生成时记录的米粒中心，不必经过教师。
把大模型的输出当作参考答案，是核对过之后才敢用的。
拿它与人工标注逐点匹配，教师给出的点在 D3 上准确率 0.993、召回率 1.000、计数误差 0.10 粒，
在 D1 上依次为 0.957、0.967 与 3.29 粒。这两组数据上的伪标签足够干净。
人工标注全程不参与训练，只用于测试。学生的成绩因此仍是拿人工标注量出来的。
图 \ref{fig:teacherlabel} 是其中一张照片的逐点对照，漏掉的和多给的都是少数。

\begin{cfigure}
  \centering
  \includegraphics[width=\textwidth]{teacher_labels.png}
  \caption{SAM 3 给出的点与人工标注的对照。绿圈为两者配对上的米粒，
  橙圈为 SAM 3 漏掉的，蓝方块为它在同一粒米上多给的点。}
  \label{fig:teacherlabel}
\end{cfigure}

\textbf{学生模型学习的内容：} 学的是密度图，
整张图的积分即为粒数。这样安排有两点考虑。其一，米粒只有 4 像素宽，要求网络给出清晰的实例边界并不现实。
让它在米粒处“鼓起一个包”却是可行的。
其二，损失函数不依赖于把相邻两粒分开，而那恰是几何方法最吃力之处。
高斯的宽度按每张图自身的米粒间距确定，因此 4 像素与 25 像素的米粒在网络看来尺度相当。
从教师的掩膜到学生的预测，整条流程见图 \ref{fig:studentpipe}。

\begin{cfigure}
  \centering
  \includegraphics[width=\textwidth]{student_pipeline.png}
  \caption{方法二的流程}
  \label{fig:studentpipe}
\end{cfigure}

\textbf{神经网络与训练：} 三层 U-Net，27 万参数，权重文件 1 MB。
按 7:1:2 划分训练、验证与测试集，划分以源图为单位。在 RTX 4080 上训练 6000 步，约 6 分钟。

\begin{ctable}
  \centering
  \caption{学生模型与本文方法、教师在同一批测试图上的平均绝对误差，单位为粒。}
  \label{tab:student}
  \tabfit{\begin{tabular}{lrrrrr}
\toprule
数据集 & 张数 & 直接数连通域 & 方法一 & 方法二 学生模型 & SAM 3 教师 \\
\midrule
D1 真实照片 & 10 & 6.20 & \textbf{2.00} & 3.91 & 2.40 \\
D2 合成粘连 & 8 & 24.38 & \textbf{0.25} & 1.29 & 1.00 \\
D3 低分辨率 & 144 & 3.23 & \textbf{0.38} & 0.42 & 0.06 \\
\midrule
合计 & 162 & 4.46 & \textbf{0.48} & 0.68 & 0.25 \\
\bottomrule
\end{tabular}}
\end{ctable}

\subsection{结果}

结果见表 \ref{tab:student}。为便于比较，表中一并列出了不作任何粘连处理时
直接统计连通域的结果 4.46 粒，四者测的是同一批图。
学生模型的总体误差为 0.68 粒，优于方法一的 1.86 粒。
在训练数据最多的 D3 上差距最明显，为 0.42 对 1.76 粒。
D1 上亦更好，为 3.91 对 4.60 粒。D2 上则是方法一更好，为 0.25 对 1.29 粒。
它仍不及教师的 0.25 粒，这符合预期，蒸馏一般不会超过教师。
但它只有 27 万参数、1 MB 权重，约为 SAM 3 的三千分之一。
每张图耗时 2 至 25 ms，比 SAM 3 快 8 到 100 倍，见表 \ref{tab:cost}。
更要紧的是不必挑选提示词。SAM 3 换一个提示词就可能一粒也检不出，
学生则把选定的那套行为固化了下来。

图 \ref{fig:studentwin} 是 D3 上的一个例子，也正是第 \ref{sec:student} 节开头那类失败。
画面左侧的木框亮边被方法一当成一处粘连块，按面积折算成 8 粒，全图数出 24 粒，真值为 14 粒。
学生模型的密度图上，木框那一带没有鼓起任何包，积分得 14.2 粒。
它没有用到更多的形状判据，只是从教师那里学到了米粒该是什么样子。

\begin{cfigure}
  \centering
  \includegraphics[width=\textwidth]{student_vs_ours.png}
  \caption{同一张低分辨率图上两种方法的对照。中图里左侧的虚线框是方法一把木框亮边
  当作粘连块后按面积折算的 8 粒，右图为学生模型预测的密度图。}
  \label{fig:studentwin}
\end{cfigure}


\section{讨论}



\subsection{一些不足}
\label{sec:limit}



\textbf{画面里有比米更整齐的东西时，会把它当成米。}
我另找过一组带直尺的照片\cite{ds_ruler}，想补充粘连数据。
结果在不少图上，程序选出的前景是直尺的刻度，而不是米粒，见图 \ref{fig:scope}。
原因在于式 \eqref{eq:score} 挑二值化方案时，看的是画面里有多少块大小一致的小东西。
刻度线大小整齐、对比强烈，比颜色浅、大小不一的米粒更符合这个标准。
换句话说，方法一只认形状，不认得什么是米。
这组照片里还有不少是几个品种的米放在一起对比，粒径相差很大，
不符合“同一张图里米粒大小相近”这一前提，所以最终没有用作测试集。

\begin{cfigure}
  \includegraphics[width=\textwidth]{scope_failure.png}
  \caption{方法被直尺刻度误导的两个例子，每例左为原图、右为选出的前景。}
  \label{fig:scope}
\end{cfigure}

\textbf{图像模糊时，可能选错二值化方案。}
第 \ref{sec:robust} 节里出现崩溃的几张图，都是模糊造成的。
模糊之后，饱和度通道被阈值切成大量小碎点。
碎点数目多，式 \eqref{eq:score} 的得分反而略高于正确的方案，程序便选错了，一张图数出上千粒。
这个问题目前还没有解决。

\textbf{硬币有浮雕时，可能被数成几粒米。}
剔除硬币的判据要求硬币先被二值化成完整的一块。
硬币表面有浮雕时，二值化后它会碎成几块，每块都和米粒差不多大，于是被当成米计了数。
附录 \ref{app:renders} 中就有这样一例。





\section{结论}

本文实现了一种基于注水分割的粘连米粒计数方法。它先由图像自身量出单粒米的尺度，
据此判断粘连、控制种子数量，最后按面积与形状复核切分结果。
三组数据上的平均绝对误差分别为 4.84、0.25 与 1.90 粒，优于五种对比方法。

做完这份作业，体会最深的一点是，这个问题没有想象中这么简单，我看到问题的第一眼，认为只要对二值化图像稍加处理即可。但遇到了挫折，于是我去网上搜集论文，看看有没有好的方法可以参考。受到了分水岭算法的启发，针对这个问题具体化，于是有了方法一。对于方法二，是一种偷懒的黑箱做法，效果好但不可解释。我做大创项目时使用过的SAM3分割模型，于是有了蒸馏学生模型的想法。

我将在文末附上更多的效果图和参考文献，头一次用Latex写报告，不太熟练，文中若有不恰当之处，还请老师批评指正。
